% Section 7: Open problems
\subsection{Open problems}
\label{sec:forecasting_rl_open}

Four problems deserve attention from researchers who already understand
econometric identification. First, action-conditional foundation forecasters do
not yet exist. The closest existing work is \citet{Ada2025}, which uses
Lag-Llama as a passive observation-process prior inside an offline RL loop, but
not as an action-conditional dynamics model. Building a TSFM whose pretraining
objective is conditional on a notion of intervention is the natural next step
and is the right architectural problem for a foundation-model team working with
an applied economist. Second, identification under TSFM-induced counterfactuals
remains unresolved. The Buesing-Oberst SCM programme assumes a correctly
specified causal graph, and the Lucas critique applies once the agent's policy
enters the data-generating process. Adding instrumental variable restrictions,
along the lines of the recent econometric literature on machine-learning
instruments, is a credible direction. Third, online decision-focused regret
bounds outside the convex polyhedral setting are at the research frontier
\citep{Capitaine2025}. Non-stationary economic environments will require dynamic
regret guarantees that are robust to path variation. Fourth, the
\citet{PerezDiaz2025} form-and-function critique calls for pretraining objectives
that produce joint trajectory ensembles rather than marginal point or parametric
forecasts. Path-dependent operational queries, including threshold-crossing
inventory rules and worst-case execution costs, are not answerable from
marginally-calibrated forecasts no matter how accurate they are.
